\appendices
\section{Archival Classification Summary}\label{app:archive}
Table~\ref{tab:architectures} and Fig.~\ref{fig:endpoints} retain the complete original configuration summary. Ordering is alphabetical. These pooled values document endpoint differences; the fold-specific results in the main text remain the primary interpretation.
\begin{table}[t]
\caption{Archival classification summary in alphabetical order. Each row averages nine runs (three grouped folds and three seeds). These internal validation summaries are retained for completeness; fold-specific analysis is primary.}
\label{tab:architectures}
\centering
\small
\setlength{\tabcolsep}{4pt}
\begin{tabular}{lrr}
\toprule
Architecture & Selected & Final \\
\midrule
CLIP ViT-B/16 & 0.95076 & 0.78897 \\
CoAtNet-0 & 0.80537 & 0.72041 \\
ConvNeXtV2-T & 0.90891 & 0.90238 \\
DeiT-III-S & 0.74726 & 0.72855 \\
DenseNet-121 & 0.89660 & 0.83852 \\
DINOv2-B & 0.91277 & 0.88380 \\
DINOv2-S & 0.87346 & 0.78842 \\
EfficientNetV2-S & 0.97615 & 0.70874 \\
MaxViT-T & 0.92189 & 0.85982 \\
MobileNetV4 & 0.99746 & 0.91310 \\
RegNetY-016 & 0.93495 & 0.74087 \\
ResNet-50 & 0.93651 & 0.82922 \\
ResNeXt-50 & 0.92260 & 0.76566 \\
Swin-S & 0.95259 & 0.85744 \\
Swin-T & 0.97673 & 0.83478 \\
VGG16-BN & 0.89194 & 0.45870 \\
ViT-S & 0.93235 & 0.75858 \\
\bottomrule
\end{tabular}
\end{table}

\begin{figure}[t]
\centering
\includegraphics[width=\columnwidth]{endpoints.pdf}
\caption{Validation-selected and fixed-final macro-F1 means for all 17 retained configurations. Each point averages nine runs on the same historical folds; connecting segments highlight endpoint differences and are not uncertainty intervals. Configurations are alphabetical, not performance-ranked.}
\label{fig:endpoints}
\end{figure}


\section{Confidence Diagnostics}\label{app:confidence}
The previously executed confidence analysis averages the saved probabilities from the three explanation-selected architectures and their three seeds within each fold. Each fold's held-out photographs are stratified into three disjoint image subsets for temperature fitting, conformal calibration, and final diagnostics, using the recorded seed $100+\mathrm{fold}$. This subdivision is image-level, not a new tyre-disjoint split. Temperature is fitted by minimising negative log-likelihood on the first subset. Expected calibration error uses 15 equal-width confidence bins. Unlike the standalone ordinal decision in~\eqref{eq:ordinal}, this saved ensemble diagnostic explicitly uses the argmax of its averaged class probabilities.

Conformal nonconformity is $1-p_y$ on the separate calibration subset. The implementation uses the higher empirical quantile at $\min\{1,\lceil(n_{cal}+1)0.9\rceil/n_{cal}\}$ and includes test class $k$ when $1-p_k$ does not exceed that threshold. Coverage and mean set size are then measured on the third subset. The recorded abstention rate counts sets larger than one, leaving empty sets outside that field. We preserve this implementation detail rather than interpreting the field as a complete reject-option metric.


For CORAL outputs $q_0=\sigma(z_0)$ and $q_1=\sigma(z_1)$, the saved implementation constructs $r=(1-q_0,q_0-q_1,q_1)$ and normalises $p_k=\max(r_k,10^{-8})/\sum_j\max(r_j,10^{-8})$. Clamping guarantees a valid probability vector even if cumulative outputs are out of order; it does not guarantee calibrated probabilities. Averaging these vectors and taking argmax can differ from counting thresholds in~\eqref{eq:ordinal}. The uncertainty analysis is consequently a separate decision pipeline, not interchangeable with the reported ordinal classifier.

\subsection{Recorded Outcomes}
Temperature scaling leaves the recorded macro-F1 unchanged and reduces confidence-error metrics on the retained calibration-test subsets (Table~\ref{tab:calibration}). On fold 1, expected calibration error decreases from 0.18955 to 0.07377, and negative log-likelihood from 0.35630 to 0.16226, with 41 test images. Extremely small calibrated errors on the other folds are not evidence of deployment-grade certainty.

At nominal 90\% coverage, the observed conformal coverages are 0.86885, 0.97561, and 0.93939 on folds 0, 1, and 2, respectively (Table~\ref{tab:conformal}). Thus only fold 0 is below nominal in these observed subsets. Mean set sizes below one on folds 0 and 2 imply empty sets. The recorded abstention field does not count all empty-set events and cannot be presented as a complete rejection policy. Small subset sizes, dependent acquisition, and reuse of internal evidence constrain generalisation of these diagnostics.
\begin{table}[t]
\caption{Confidence diagnostics on held-out calibration-test subsets. ECE is the recorded expected calibration error; values near zero on internal folds do not imply external reliability.}
\label{tab:calibration}
\centering
\small
\setlength{\tabcolsep}{4pt}
\begin{tabular}{llrrrr}
\toprule
Fold & Scores & $n$ & F1 & ECE & NLL \\
\midrule
0 & Raw & 61 & 1.0000 & 0.2115 & 0.2425 \\
0 & Temp. & 61 & 1.0000 & 0.0000 & 0.0000 \\
1 & Raw & 41 & 0.8030 & 0.1896 & 0.3563 \\
1 & Temp. & 41 & 0.8030 & 0.0738 & 0.1623 \\
2 & Raw & 33 & 1.0000 & 0.1293 & 0.1438 \\
2 & Temp. & 33 & 1.0000 & 0.0000 & 0.0000 \\
\bottomrule
\end{tabular}
\end{table}

\begin{table}[t]
\caption{Prediction-set diagnostics at nominal 90\% coverage. The recorded abstention field does not count every empty-set event.}
\label{tab:conformal}
\centering
\small
\setlength{\tabcolsep}{4pt}
\begin{tabular}{rrrrrr}
\toprule
Fold & $n_{cal}$ & $n_{test}$ & Coverage & Set size & Abstain \\
\midrule
0 & 62 & 61 & 0.8689 & 0.8689 & 0.0000 \\
1 & 43 & 41 & 0.9756 & 1.0976 & 0.0976 \\
2 & 35 & 33 & 0.9394 & 0.9394 & 0.0000 \\
\bottomrule
\end{tabular}
\end{table}


